\documentclass[../main.tex]{subfiles}
\begin{document}
%==========================================TIÊU ĐỀ TẬP========================================
\begin{table}[h]
    \begin{adjustwidth}{-1cm}{}
        \begin{tabular}{>{\centering\arraybackslash}p{6cm}|>{\raggedright\arraybackslash}p{12.5cm}}
            \multirow{5}{6cm}
            % Cột bên trái: ÉPISODE và số tập
            {\\[-40pt]
                \flaregothic\fontsize{20pt}{20pt}\selectfont \centering
                \color{eptype}ÉPISODE\color{black}\\
                \burbank\fontsize{72pt}{72pt}\selectfont
                \color{epnum}37\color{black}\\[6pt]
                \cafeteria\fontsize{20pt}{20pt}\selectfont\color{epnum}(3-8)\color{black}
                \\[18pt]
                \flaregothic\fontsize{14pt}{14pt}\selectfont
                \color{epattr}MISE À JOUR DU MODÈLE pour \uline{\color{Mikopi} Miko}\\
                \flaregothic\fontsize{18pt}{18pt}\selectfont
                \color{holofes}L'HOLOÊTE !\\
                \flaregothic\fontsize{18pt}{18pt}\selectfont
                \color{epattr}ÉPISODE PERDU
                \color{black}
            }
             &
            % Dòng thứ nhất: tên tập tiếng Nhật
            {\fotskip\fontsize{18pt}{18pt}\selectfont \ruby{大}{だい}\ruby{事}{じ}な\ruby{ラ}{ら}\ruby{イ}{い}\ruby{ブ}{ぶ}に\ruby{参}{さん}\ruby{加}{か}できない\ruby{ド}{ど}\ruby{ッ}{っ}\ruby{キ}{き}\ruby{リ}{り}
            }                                                                                     \\[6pt]
            % Dòng thứ hai: tên tập tiếng Anh
             & {\cafeteria\fontsize{18pt}{18pt}\selectfont Prank goes Nonstop}                                                                             \\[4pt]
            % Dòng thứ ba: tên tập tiếng Việt
             & {\vnmsans\fontsize{18pt}{18pt}\selectfont Cú lừa liên hoàn}                                                                                 \\[8pt]
            % Dòng thứ tư: Nhân vật xuất hiện
             & {\caviar{\fontsize{14pt}{14pt}\selectfont Nhân vật: \emp{kuon1} \emp{ryuuhou1} \emp{achan} \emp{sora} \emp{roboco1} \emp{miko2}}}
            \\[8pt]
            % Dòng cuối cùng: Thời gian diễn ra của tập (thường sẽ là ngày phát hành)
             & {\continuum\fontsize{16pt}{16pt}\selectfont Ngày phát hành: 19/12/2019}                                                                           \\[18pt]
        \end{tabular}
    \end{adjustwidth}
\end{table}
%=========================================TRUYỆN CHÍNH========================================
\color{black}

\txtbx[2]{{\color{vol} \uline{Tóm tắt:}} Cả văn phòng hololive náo loạn khi Tokino Sora bất ngờ mắc chứng ``mất trí nhớ'' quái đản ngay trước thềm buổi đại nhạc hội Toyosu PIT. Miko và Roboco hoảng hốt tìm đủ mọi cách từ tra nhật ký đến\dots\ cho ăn cá để cứu chữa người chị cả, mà không hề hay biết mình đang rơi vào một cái bẫy tinh vi. Trong khi đó, Ryuuhou - người đã thấu tỏ mọi chuyện - lại hăng hái ``diễn sâu'' nhất để đẩy sự hỗn loạn lên đỉnh điểm, biến buổi sáng bình thường thành một màn kịch đầy rẫy những con cá khổng lồ và những cú lừa không hồi kết.}

\begin{center}
    \textit{\uline{(Số này được thực hiện nhằm quảng bá cho buổi concert Toyosu PIT diễn ra vào ngày 24/1/2020.)}}
\end{center}

Một buổi sáng mùa đông nắng hanh vàng tại văn phòng \hll. Kuon đang ngồi cùng Miko, Roboco và cô em gái Ryuuhou để rà soát lại danh sách hạng mục cho buổi concert sắp tới. Miko hôm nay trông vô cùng rạng rỡ trong bộ trang phục mới - một thiết kế hiện đại lấy cảm hứng từ đồ Vu nữ truyền thống nhưng hở lưng táo bạo, kết hợp với đôi vớ không đối xứng đầy phong cách.

Bộ trang phục của em ấy về cơ bản là giống với bộ trang phục của một người vu nữ, nhưng mang hơi hướng phi truyền thống hơn: áo kimono hở lưng cả ở lưng và ở hông không tay, có bèo nhún với cổ áo sơ mi thông thường và tua rua (dây buộc sau lưng) màu đỏ. Cô có hai tay áo rộng màu trắng được tách rời, váy ngắn hakama màu đỏ, dây đai đùi màu đỏ ở một chân và dây garter ở chân còn lại. Cô đi đôi tabi không đều (một chiếc tất trắng và một chiếc quần tât cao đến đùi có viền ren trắng), đi cùng với dép geta hồng. Miko giờ đã có tóc dài đến lưng, với một búi tóc được buộc lệch sang bên trái bằng chiếc chuông vàng cột với nơ hồng.

\model{24}{l}{6cm}{miko2}

``Nào, báo cáo tình hình cho anh xem nào,'' Kuon gác chân lên bàn, vẻ mặt thư thái.

Roboco nhanh nhảu lật quyển sổ tay, giọng hào hứng: ``Đại nhạc hội tại Toyosu PIT sẽ diễn ra vào cuối tuần này rồi ạ! Mọi thứ đang vào guồng quay chóng mặt luôn!''

``Rồi. Những người khác thì sao? Miko-chan, em có ý kiến gì không?''

Miko gật đầu lia lịa, tay chỉnh lại chiếc chuông vàng trên búi tóc: ``Đúng đó ne! Miko đã chuẩn bị xong xuôi cả danh sách bài hát rồi!''

Kuon liếc nhìn sang Ryuuhou, lúc này đang thản nhiên gặm một chiếc taiyaki: ``Còn em thì sao, Ryuuhou? Đừng bảo là em chỉ đến đây để ăn chực đấy nhé.''

Cô bé nheo đôi mắt hai màu đỏ - xanh lam, nở một nụ cười đầy ẩn ý: ``Em á? Em đang giám sát tiến độ để đảm bảo không có ai\dots\ quên mất việc quan trọng nhất thôi mà.''

``Giám sát cái con khỉ mốc. Có mà em kiếm cớ để đi chơi thì có.''

Vừa dứt lời, một giọng nói đầy ngạc nhiên vang lên từ phía cửa: ``\textbf{Thật sao!?}''

Cả bốn người đồng loạt quay lại.  Sora đang đứng đó, khuôn mục vốn dĩ dịu dàng nay lại hiện rõ vẻ ngơ ngác đến lạ lùng. ``Mình không hề biết là chúng ta sắp có đại nhạc hội đấy\dots''

Ánh sáng trong mắt nàng Thần tượng dường như đã biến mất hoàn toàn, thay vào đó là một cái nhìn xa xăm, vô định.

\img{3-8a}

``Sora-chan? Em đến từ khi nào vậy?'' Kuon hỏi, giọng điệu có chút ngờ vực.

``Em vừa mới đến thôi ạ\dots'' Sora lẩm bẩm, rồi quay sang nhìn Miko như thể nhìn một người lạ mặt. ``Cho mình hỏi\dots\ bạn là nhân viên mới của văn phòng à?''

Cả căn phòng lập tức rơi vào trạng thái đóng băng trong tích tắc. Miko há hốc mồm, chiếc chuông trên tóc khẽ kêu lên một tiếng 'keng' vì chủ nhân của nó đang run rẩy.

``Sao lại như thế được ne!?'' Miko hét lên trong sự hoảng loạn.

Roboco ôm đầu, các cảm biến trên người bắt đầu nháy đèn báo động: ``Trời ơi! Không lẽ cậu ấy bị mất trí nhớ nghiêm trọng rồi sao!?''

Ryuuhou ngay lập tức bỏ miếng bánh xuống, nhảy bật dậy với vẻ mặt kinh hoàng chưa từng thấy: ``Ôi không! Sora-san không nhận ra Miko-san sao!? Chuyện này còn tệ hơn cả tận thế nữa! Có khi nào chị ấy quên luôn cả cách thở không!?''

Kuon liếc nhìn cô em gái, thầm nghĩ trong đầu: ``\textit{Ảo tưởng cô nương.}'' Tuy nhiên, cậu vẫn phải hắng giọng lại để giữ giọng: ``Bình tĩnh nào mọi người. Có lẽ chỉ là do em ấy làm việc quá sức nên bị đãng trí nhất thời thôi.''

Sora khẽ cười, một nụ cười ngây ngô đến phát sợ: ``Đúng đó! Mình vẫn nhớ được tên mìn mà, thậm chí còn viết được nữa đây này! Xem nhé!''

Nói rồi, cô nhanh chóng lôi ra một tấm bảng trắng cùng một vật thể sắt nhọn mà cô khẳng định là `cây bút'. Ryuuhou vội vàng cổ vũ: ``Đúng rồi! Sora-senpai giỏi lắm, viết cho mọi người xem đi chị!'' Tôi nghĩ rằng em ấy đang muốn trêu chọc mọi người, dù hôm nay không phải ngày Cá tháng Tư.

\img{3-8b}

``Ờm\dots\ Sora-chan, em có chắc đó là cây bút không?'' Kuon hỏi, nhướng mày nhìn cái cờ lê trong tay cô.

Sora không đáp, cô nghiêm túc di cái cờ lê lên bảng, phát ra những tiếng 'két két' chói tai. Tuy nhiên, chẳng có lấy một nét mực nào hiện ra. Sora nghiêng đầu, bĩu môi đầy thất vọng: ``Ủa? Sao mình không viết được bằng cái dĩa này nhỉ\dots''

``Cái đó là cái cờ lê mà! Với cả ai lại đi viết bằng dĩa bao giờ!'' Miko gần như phát điên, tay chỉ trỏ loạn xạ.

Ryuuhou ôm mặt, rên rỉ: ``Thôi xong rồi! Chị ấy đã mất khả năng phân biệt dụng cụ cơ bản! Cứ đà này, chị ấy sẽ dùng bàn chải bồn cầu để đánh răng mất thôi!''

Roboco chống hông, vẻ mặt nghiêm trọng: ``Nếu cậu ấy còn mất trí thêm chút nữa, cậu ấy sẽ quên luôn cả `đạo đức nghề nghiệp' mất! Thần tượng mà quên đạo đức nghề nghiệp là thảm họa đấy!''

Kuon thở dài, bước đến trước mặt Sora, đặt tay lên vai cô: ``Được rồi, Sora-chan. Nghe anh hỏi này: Chữ Hán đầu tiên trong tên anh (久) có nghĩa là gì?''

Sora trầm tư suy nghĩ hồi lâu, rồi tươi cười trả lời: ``Là 'Xa rời' ạ!''

Tôi nhíu mày, thầm nghĩ có lẽ đây chỉ là khởi đầu. Tôi tiếp tục đặt câu hỏi khác: ``OK\dots\ Thế còn tên của em, 'Sora' viết theo Hán tự nghĩa là gì?''

``Đất ạ!'' Sora đáp đầy tự tin.

Đến đây thì tôi hoàn toàn bó tay, chỉ biết tặc lưỡi, lắc đầu cười ngao ngán: ``Miko-chan và Roboco-chan nói đúng. Em ấy mất trí nhớ luôn rồi.''

Miko nghe xong thì ngã quỵ xuống sàn: ``Không ổn rồi! Chị cả của chúng ta hỏng thật rồi ne! Mau gọi bác sĩ, thầy cúng, hay bất cứ ai có thể hồi phục bộ não này đi!''

``Để mình thử dùng trí tuệ nhân tạo tra cứu nhật ký hệ thống xem sao,'' Roboco tuyên bố, rồi ngồi xuống, toàn thân phát ra ánh sáng xanh lấp lánh như một chiếc đèn disco.

Miko ngay lập tức reo lên đầy phấn khích: ``Quả đúng là rô-bốt chất lượng cao của chúng ta mà, ne!'' Tôi thì nhếch mép, đáp lại bằng giọng mỉa mai mỗi lần nhắc đến em ấy: ``Ờ. Chất lượng cao đến mức mà phá luôn cả cái công ty bữa trước là đủ hiểu nó chất lượng thế nào rồi đấy.''

Robo-PON lờ đi câu nói giễu cợt của tôi, bình thản ngồi xuống. Em ấy bắt đầu suy nghĩ cách giải quyết, và trong quá trình đó, cả cơ thể bỗng phát sáng lấp lánh.

\img{3-8c}

``Chỉ là tra cứu thôi mà cần phải rực rỡ thế sao ne?'' Miko lẩm bẩm trong nước mắt. Ryuuhou thì gật gù tán thưởng: ``Hiện đại thì nó phải hại điện thế chứ chị!'' 

Sau một hồi ``xử lý dữ liệu'', Roboco rút ra một mảnh giấy, vẻ mặt đầy tự hào. Tuy nhiên, nội dung trên đó lại khiến cả căn phòng chết lặng thêm lần nữa: \textit{``Năm Tý - dầu gội ngon quá.''}

\img{3-8d}

Kuon không nhịn được mà bật cười ha ha, chỉ tay vào nàng Người máy: ``Đúng là 'chất thải' cao thật! Roboco-chan, em ăn cả dầu gội đầu à?''

``Đó là cảm nhận về hương vị thôi mà! Đừng có trêu em nữa!'' Roboco đỏ bừng mặt, giậm chân bành bạch.

Trong lúc hỗn loạn, Sora lại tiếp tục làm mọi người sốc bằng cách trình diễn kỹ năng `đi bộ bằng hai chân xen kẽ' như một đứa trẻ mới tập đi, hoặc nhẹ hơn, như một cô bé 18 tuổi lần đầu được tập đi sau khi được chữa liệt người. Đôi chân cô vung lên cao quá mức, dáng đi cứng ngắc như một con bù nhìn gỗ bị lỗi phần mềm. ``Mọi người ơi, xem này! Mình có thể đi về phía trước bằng cách bước chân trái và chân phải xen kẽ nhau này!''

``Cái đó thì ai chẳng làm được ne!'' Miko ôm đầu gào thét. Cô nàng ôm đầu quay qua quay lại như thể tìm kiếm giải pháp trong vô vọng, mắt của em ấy thậm chí còn xoay tít thò lò: ``Bà ấy không thể cứ như thế này mãi được ne! Làm sao bây giờ đây!?''

Ryuuhou đế thêm vào: ``Nhìn kìa! Chị ấy đang tiến hóa ngược thành động vật bậc thấp rồi! Onii-san, hay là mình lấy cái gì đó gõ vào đầu chị ấy để 'reset' lại hệ thống đi?''

``Em định giết người à!?'' nàng Vu nữ hốt hoảng ngăn cản.

Đúng lúc này, Roboco bất ngờ lôi từ sau lưng ra một con cá nhỏ vẫn còn đang giãy giụa: ``Mình nghe nói ăn cá sẽ giúp chúng ta thông minh hơn đấy. Thử xem.''

Miko nhìn con cá với ánh mắt chết lặng: ``Tôi nghĩ bà cần ăn con cá đó hơn cậu ấy đấy ne.''

Tôi góp thêm một câu, vẫn giữ nguyên vẻ trêu chọc: ``Và cả ít việt quất nữa. Ăn cho khai thông trí thức ra đi chứ em toàn phát biểu mấy cái đẩu đâu ấy.''

\img{3-8f}

Trong khi ba cô gái còn đang chí chóe quanh con cá, Kuon và Sora lén liếc nhìn nhau, rồi liếc về phía góc phòng nơi A-chan đang bí mật đặt máy quay. Ryuuhou cũng nháy mắt với Sora, môi khẽ nhếch lên một nụ cười đắc ý.

\begin{center}
    \textit{\uline{30 phút trước tại Phòng điều hành.}}
\end{center}

Kuon, Sora, A-chan và Ryuuhou đang ngồi quanh bàn họp, bàn bạc về kế hoạch quảng bá cho buổi concert sắp tới. A-chan, với kinh nghiệm dày dặn trong việc tạo ra những nội dung viral, đang cố gắng nghĩ ra một ý tưởng độc đáo để thu hút sự chú ý của khán giả.

``Chúng ta cần một ý tưởng nổ não để quảng bá cho Toyosu PIT,'' A-chan gõ bút xuống bàn. ``Có ai nghĩ được cách nào để quảng bá cho buổi còn-sớt tiếp theo cho công ty không?''

Tôi giơ tay lên, và khi được mời phát biểu, tôi nói ra ý tưởng của mình: ``Sao chúng ta không dùng số tiếp theo của \textit{holo no Graffiti} để quảng bá nhỉ? Giống như cách mà mọi người từng làm với Comiket năm ngoái ấy?''

Cô nàng Đeo kính gật gù, đồng tình với quan điểm: ``Cậu cũng có ý tưởng giống tôi đấy. Nhưng mà chúng ta sẽ làm kịch bản theo cách nào?''

``Hay là chơi khăm đi?'' Sora hào hứng đề xuất. ``Em sẽ giả vờ bị mất trí nhớ, xem phản ứng của Miko và Roboco thế nào!''

Kuon gật đầu tán thành: ``Nghe được đấy. Anh sẽ đóng vai người quan sát để dẫn dắt câu chuyện.''

Ryuuhou giơ tay xung phong: ``Cho em tham gia nữa! Em sẽ đóng vai người 'hùa theo' để dọa cho hai chị ấy sợ phát khiếp luôn! Em hứa sẽ không cười lúc đang diễn đâu!''

A-chan cười rộ lên: ``Được đấy! Vậy chốt kịch bản 'Sora mất trí nhớ' nhé!''

Thế là, từ việc phân vai đến tình huống cụ thể, tất cả đều được lên kế hoạch cẩn thận. Mọi thứ diễn ra ở trên chính là kết quả của buổi họp này: một trò chơi khăm mà nàng Thần tượng mạnh dạn vào vai ``kẻ điên mất trí,'' còn Kuon thì đóng vai ``nạn nhân ngây thơ,'' một ``vai diễn bất đắc dĩ'' mà cậu luôn đóng mỗi khi đi làm.

\ct

\begin{center}
    \uline{Hiện tại.}
\end{center}

Quay lại với thực tại, Miko và Roboco vẫn đang cuống cuồng tìm cách ``cứu'' Sora. Kuon tranh thủ ghé sát vào máy quay ẩn, thì thầm: ``\hl{Đến giờ hai đứa nó vẫn chưa biết gì cả. Kế hoạch của A-san thành công mỹ mãn rồi.}''

Sora cũng ghé đầu vào, cười khúc khích: ``\hl{Họ thực sự rất quan tâm đến em đấy! Nhưng mà vui quá, em phải diễn tiếp thôi!}''

Ryuuhou thình lình xuất hiện sau lưng hai người, cũng ghé mắt vào ống kính: ``\hl{Sắp đến đoạn cao trào rồi! Mọi người chuẩn bị xem A-chan xuất hiện nhé!}''

Cậu Tổng đứng thẳng lên, hắng giọng để tiếp tục vào vai. Tôi tiến tới hai cô nàng đang hoảng loạn kia: ``Nào nào, Miko-chan, Roboco-chan! Giờ ta phải tìm cách giúp em ấy ngay!''

Sora lập tức quay lại với biểu cảm ``mất trí nhớ'' ngây ngô: ``Ủa, giúp em gì cơ? À, đúng rồi! Em vừa nhớ ra mình có thể\dots\ vừa đi vừa nhảy một chân!''

Câu nói của nàng khiến cả Roboco và Miko sững sờ trong vài giây, rồi lại tiếp tục lao vào tranh luận xem phải xử lý tình huống này như thế nào. Trong lòng cậu thầm nghĩ: ``\textit{Má, có khi đóng kịch thì họ nhập vai còn giỏi hơn cả Sora-chan nữa.}''

Trong lúc đó, nàng Người máy bỗng lấy ra một con cá khổng lồ khác, to hơn hẳn con trước đó: ``Còn con cá này thì sao? Nó còn to hơn nữa đấy!''

Miko hốt hoảng thốt lên: ``Bà kiếm con cá đó ở đâu ra vậy ne!? Bắt ở dưới biển à?''

Trong khi hai người tranh cãi chí chóe, cả ba người ngoài đứng lùi về sau, cười thầm và nghĩ bụng: ``\textit{Tiếp theo sẽ là A-chan đi vào và kết thúc mọi chuyện.}''

Nhưng sau một lúc, cậu bắt đầu nhận thấy em ấy có chút sốt ruột. Sora hướng về phía cửa, ánh mắt giương ra kiểu: ``\textit{Sao A-san đến lâu vậy nhỉ?}''

Đúng lúc đó, cửa phòng bật mở. A-chan lao vào như một cơn lốc, khuôn mặt lộ rõ vẻ 'hớt hải' đã được tập dượt kỹ lưỡng: ``Chuyện gì thế này!? Chị nghe nói Sora gặp rắc rối lớn!''

Miko nhìn thấy chị ấy lao vào, như thể bắt được vàng, vội vã thông báo: ``A-chan! Chị thấy đấy, Sora-chan đã\dots''

\begin{figure}[H]
    \centering
    \begin{subfigure}[t]{0.45\textwidth}
        \centering
        \includegraphics[width=8cm]{../../../../Image/Book3/3-8h1.png}
        \caption{A-chan ập vào, đang định kết thúc vở kịch với một Miko đang hoảng loạn\dots}
    \end{subfigure}
    \quad
    \begin{subfigure}[t]{0.45\textwidth}
        \centering
        \includegraphics[width=8cm]{../../../../Image/Book3/3-8h2.png}
        \caption{\dots nhưng rồi cô bị con cá khổng lồ mà Roboco trượt tay đập trúng mặt.}
    \end{subfigure}
\end{figure}

Nhưng khi cô nàng còn chưa phân trần xong thì đột nhiên\dots\ *BỐP*! Nguyên một con cá siêu to đáp trọn vào mặt A-chan, khiến chị ấy ngã ngửa xuống đất.

\img{3-8i}

``Ối chết! Tay em bị trơn!'' Roboco cuống quít đỡ A-chan dậy.

Kuon mắng át đi: ``Tên hậu đậu này! Làm hỏng hết việc rồi!''

A-chan lờ đờ ngồi dậy, đưa tay ôm đầu, ánh mắt mơ màng hỏi một câu chí mạng: ``Đây là đâu\dots\ Tôi là ai?''

``\textbf{Trời ơi! Lại thêm một người nữa bị lây bệnh mất trí nhớ rồi ne!!}'' nàng Vu nữ hét lên trong tuyệt vọng, tiếng hét vang vọng khắp cả văn phòng.

Sora thì vẫn đứng thăng bằng, nhưng khuôn mặt nàng Thần tượng hiện rõ vẻ hoang mang. Nhưng Kuon vẫn giữ vẻ điềm tĩnh, vì trong lòng đã rõ: ``\textit{Chắc chắn chị ấy đang chơi khăm tiếp. Đúng là `gừng càng già càng cay' mà!}''

Ryuuhou lăn ra sàn cười sặc sụa (vẫn cố giả vờ là đang khóc): ``Thôi xong rồi, A-chan cũng quên luôn cả cách trả lương cho chúng ta rồi sao!? Đáng sợ quá!''

\begin{figure}[H]
    \centering
    \begin{subfigure}[t]{0.45\textwidth}
        \centering
        \includegraphics[width=8cm]{../../../../Image/Book3/3-8j1.png}
    \end{subfigure}
    \quad
    \begin{subfigure}[t]{0.45\textwidth}
        \centering
        \includegraphics[width=8cm]{../../../../Image/Book3/3-8j2.png}
    \end{subfigure}
\end{figure}

Lúc này, Miko chợt nhận thấy A-chan đang cầm một tấm bảng giấu sau lưng. Cô lao tới giật lấy, và thấy trên đó ghi dòng chữ: ``CÚ LỪA ĐẤY!''

Dù đã rõ như ban ngày như thế, nàng Người máy lẩm bẩm một câu ngờ nghệch, hoàn toàn không biết chuyện gì đang xảy ra: ``Không có con giáp nào là Cá sao?'' Cậu Tổng bật cười, trả lời ngắn gọn: ``Không có đâu mẹ trẻ ơi.''

Cả nhóm còn chưa kịp thảo luận thêm, thì bất ngờ, từ trên trần, một sợi dây được thả xuống. Roboco-chan kéo sợi dây trong tò mò, và ngay lập tức, một quả pháo nhỏ bắn ra, kèm theo một mảnh giấy tung bay trong không trung. Dòng chữ trên đó ghi rõ ràng: \textit{``Bạn đã bị chơi khăm! Nhớ mua vé đại nhạc hội Toyosu PIT nhé!''}

\begin{figure}[H]
    \centering
    \begin{subfigure}[t]{0.45\textwidth}
        \centering
        \includegraphics[width=8cm]{../../../../Image/Book3/3-8k1.png}
    \end{subfigure}
    \quad
    \begin{subfigure}[t]{0.45\textwidth}
        \centering
        \includegraphics[width=8cm]{../../../../Image/Book3/3-8k2.png}
    \end{subfigure}
\end{figure}

Miko nhìn mảnh giấy, rồi hét lên đầy bức xúc: ``\textbf{Vậy thì cái bảng để làm gì chứ!?}'' Ryuuhou nhún vai, khẽ bật cười: ``Đề phòng ai đó vẫn chưa hiểu thôi mà, Miko-san!''

Cả căn phòng chìm trong tiếng cười rộn rã, ngoại trừ Miko và Roboco vẫn còn đứng đực ra như hai bức tượng vì không thể tin nổi mình vừa bị 'ăn thịt lừa' một cách ngoạn mục như vậy.

%==========================================TRUYỆN PHỤ=========================================
\omk

Sau khi những tràng cười lắng xuống, cả nhóm quây quần lại quanh bàn trà. Sora và A-chan không bỏ lỡ cơ hội, bắt đầu tung ra một loạt tờ rơi quảng cáo cho buổi concert sắp tới.

Khi mọi thứ sáng tỏ, hai khuôn mặt trái ngược nhau hiện rõ: nàng Vu nữ thở phào nhẹ nhõm, trong khi nàng Robo-Ngốc thì vẫn ngơ ngác như chưa hiểu hết chuyện ma xơi quỷ hờn gì.

Miko vẫn còn phồng má tức tối, nhìn Kuon trân trân: ``Anh cũng tham gia vào chuyện này sao ne!? Đồ xấu tính! Miko đã thực sự lo lắng cho Sora-chan lắm đó!''

Kuon nhún vai, cười đắc thắng: ``Thì anh là quản lý mà, phải biết cách tạo kịch tính chứ. Với lại em nhìn xem, phản ứng của em lúc nãy đáng giá triệu view đấy.''

Ryuuhou vỗ vai Miko, vẻ mặt đầy đắc ý: ``Em cũng biết hết rồi đó chị Miko! Lúc em bảo chị Sora sắp quên cách thở, chị tin sái cổ luôn kìa! Buồn cười chết đi được!''

``Ryuuhou! Ngay cả em cũng lừa chị sao!?'' Miko hét lên, định nhào tới cù lét cô bé nhưng Ryuuhou đã nhanh chân nấp sau lưng Kuon.

Roboco thì vẫn đang chăm chú nhìn con cá khổng lồ trên bàn, thắc mắc: ``Thế rốt cuộc là không có con giáp nào là con Cá sao ạ? Em thấy nó to thế này cơ mà\dots''

Kuon cốc nhẹ vào đầu nàng Người máy: ``Ăn việt quất vào cho tỉnh táo đi cô nương. Thông minh kiểu này thì hỏng hết cả hệ thống Staff của anh mất thôi!''

Ở góc phòng, Sora và A-chan đang xem lại đoạn băng ghi hình, thỉnh thoảng lại phá lên cười khi thấy cảnh A-chan bị 'ăn cá'.

``Chị nghĩ buổi concert lần này sẽ thành công rực rỡ đấy,'' A-chan gật gù hài lòng.

Sora mỉm cười, ánh mắt tràn đầy hy vọng: ``Vâng! Nhờ sự ủng hộ của mọi người và 'cú lừa' ngoạn mục này, chắc chắn chúng ta sẽ có một đêm diễn không thể nào quên!''

Kuon nhìn những cô gái đang vui đùa bên nhau, thầm nhủ: ``\textit{Dù có hơi điên rồ và quái quặc, nhưng chính sự náo nhiệt này mới là điều khiến văn phòng hololive trở thành nhà.}''

\end{document}
